Two bounded searches revisited the existing split-macro approach.
The first exposed each D/H/A copy as a shared Push prefix followed by its
loader; the second also exposed the closed phases of Push and constructor
arithmetic, with exact existing reader pairs available without a new
control. These are credited extensions/replays of an existing strategy,
not a new method. Each run had a 240-second budget.
\begin{center}\small
\begin{tabular}{@{}lrrr@{}}
\toprule
Search & Proposals & Accepted slots & Best fitness \\
\midrule
Copy fragments & 3\,966\,253 & 708\,450 & 194 \\
Closed phases & 4\,215\,199 & 222\,291 & 194 \\
\bottomrule
\end{tabular}
\end{center}
Across 8\,181\,452 proposal slots neither run reached the prospective
193-state gate. An independent resolver checked all 60 retained endpoint
grammars and proved their complete resolved 3673-instruction graphs equal
to the baseline. This verifies source binding only; the 60 endpoints are
not all newly certified literal machines.

Six further logical-emitter searches tried genuine antecedent exchanges
in the two membership-rewrite heads, parameter permutations, and alternative
separation matrices. The matrix words received independent full-emitter
free-term checks. The searches completed 30\,600\,000 proposed mutations;
the best allocator count was 196. The 197--200 results for the rewrite-head
searches are optimistic pair-matching scores; only the two matrix endpoints
at 196 also passed the evaluator's simultaneous-fold compatibility check.
None underwent full physical acceptance, and none is claimed as a new
machine. The frozen source, inputs, witnesses, logs, and these distinctions
are in \file{search/logic/}.

Proposal slots may repeat a tree, be no-ops, or be skipped. Neither total
counts distinct machines. All conclusions are limited to the specified
finite heuristic searches. They establish no global lower bound and do not
rule out a mathematically different route below 194.
